% Configurazione
\documentclass{article}

\usepackage{titling} % Required for inserting the subtitle
\usepackage{graphicx} % Required for inserting images
\usepackage{tabularx} % Per l'ambiente tabularx (tabelle)
\usepackage{calc} % Sempre per le tabelle
\usepackage[hidelinks]{hyperref} % Per i collegamenti ipertestuali, ad esempio sulla table of contents
\usepackage[italian]{babel} % Per la lingua italiana nelle scritte automatiche
\usepackage{xcolor} % Per colorare il testo
\usepackage{colortbl} % Per colorare le celle delle tabelle
\usepackage{lipsum} % Per generare lorem ipsum
\usepackage[normalem]{ulem} % Per sottolineare il testo
\usepackage{array} % Per la visualizzazione fluttuante di array di domande e risposte
\usepackage{ragged2e} % Pacchetto necessario per \justifying che giustifica il testo di tabelle

\newcommand{\ulhref}[2]{\href{#1}{\uline{#2}}} % Nuovo comando per sottolineare i link
\newcommand{\ulref}[1]{\uline{\ref{#1}}} % Nuovo comando per sottolineare i collegamenti a immagini
\setlength{\parindent}{0pt} % Rimuove il rientro automatico dei paragrafi

\graphicspath{ {immagini/} {../../../shared/immagini/} }


% Struttura
\begin{document}

\thispagestyle{plain} % Niente intestazione e piè di pagina


\begin{tikzpicture}[remember picture, overlay]
    % Punto di partenza al centro orizzontale nella metà superiore
    \coordinate (top_center) at ($(current page.north)!0.3!(current page.south)$);

    % UniPD: Logo e descrizione
    \node at (top_center) [anchor=north, xshift=-3cm, yshift=4.85cm] 
        {\includegraphics[width=0.15\textwidth]{Logo Universita di Padova.png}};
    \node at (top_center) [anchor=north, xshift=1.7cm, yshift=4.5cm]
        {\textcolor{red}{\textbf{Università degli Studi di Padova}}};
    \node at (top_center) [anchor=north, xshift=1.7cm, yshift=4cm]
        {\textcolor{red}{Laurea: Informatica}};
    \node at (top_center) [anchor=north, xshift=1.7cm, yshift=3.5cm]
        {\textcolor{red}{Corso: Ingegneria del Software}};
    \node at (top_center) [anchor=north, xshift=1.7cm, yshift=3cm]
        {\textcolor{red}{Anno Accademico: 2024/2025}};

    % SWEg Labs: Logo e descrizione
    \node at (top_center) [anchor=north, xshift=-2.85cm, yshift=1.5cm] 
        {\includegraphics[width=0.16\textwidth]{Logo SWEg.png}};
    \node at (top_center) [anchor=north, xshift=1.7cm, yshift=0.5cm]
        {\textbf{Gruppo: SWEg Labs}};
    \node at (top_center) [anchor=north, xshift=1.7cm, yshift=0cm]
        {Email: \textsf{gruppo.sweg@gmail.com}};
\end{tikzpicture}


\vspace{10cm}

{
\centering
\Huge\bfseries Norme di progetto\par
\vspace{0.5cm}
\Large Versione 1.0.0\par
}

\vspace{2cm}
\begin{center}
\begin{tabular}{r|l}
    \textbf{Versione} & 1.0 \\
    \hline
    \textbf{Stato} & Approvato \\
    \hline
    \textbf{Uso} & Esterno \\
    \hline
    \textbf{Approvazione} & \emph{Tutto il gruppo} \\
    \hline
    \textbf{Redazione} & Federica Bolognini \\
    \hline
    \textbf{Verifica} & ....... \\
    \hline
    \textbf{Distribuzione} & \emph{SWEg Labs} \\
    \hline
    \textbf{Destinatari} & Prof. Tullio Vardanega \\
                         & Prof. Riccardo Cardin \\
\end{tabular}
\end{center}


\newpage
\section*{Registro delle modifiche}

\begin{table}[h]
    \centering
    \begin{tabular}{|l|l|l|p{5cm}|}
        \hline
        \rowcolor[gray]{0.9}
        \textbf{Vers.} & \textbf{Data} & \textbf{Autore} & \textbf{Descrizione}\\
        \hline
        1.0 & 30-10-2024 & \emph{Tutto il gruppo} & Approvazione del documento\\
        \hline
        0.9 & 30-10-2024 & Giacomo Loat & Verifica del documento\\
        \hline
        0.8 & 29-10-2024 & Federica Bolognini & Fine stesura del documento\\
        \hline
        0.7 & 28-10-2024 & Federica Bolognini & Revisione e aggiornamento della suddivisione in fasi del progetto a seguito del confronto con il proponente\\
        \hline
        0.6 & 27-10-2024 & Federica Bolognini & Studio preliminare preventivo costi ed impegni orari\\
        \hline
        0.5 & 26-10-2024 & Federica Bolognini & Studio e suddivisione del documento in sezioni\\
        \hline
        0.1 & 21-10-2024 & Federica Bolognini & Creazione del documento\\
        \hline
    \end{tabular}
\end{table}

\newpage
\tableofcontents
\newpage

\section{Introduzione}

Esempio che usa una
\emph{parola da glossario}\textsubscript{\textit{\textbf{G}}}


\end{document}
